\subsection{A three-state stalk at a five-element middle corner}

\begin{lemma}[Proved: a small left act of the four-element rectangular band]
\label{CP2:A-three-state-P5-left-module}
Let $P_5=(\{0,1\}\times\{0,1\})^1$, with
$e_{ij}e_{k\ell}=e_{i\ell}$, and put $R=\mathbb QP_5$.
If a nonempty left $P_5$-set $S$ has at most three elements and
has no global fixed point, then $\operatorname{pd}_R\mathbb Q[S]\leq1$.
\end{lemma}
\begin{proof}
Write $K=\{e_{ij}\}$ and $U=KS$. For $u\in U$, the orbit
$Ku$ has at most two elements, one for each row. Two such orbits
are equal or disjoint: if $u=e_{ij}s$, then
$e_{k\ell}u=e_{kj}s$, and applying either row idempotent to any
point of this orbit recovers its point in that row. A singleton
orbit is fixed by every element of $K$, hence by all of $P_5$.
It is excluded by hypothesis. Since $|S|\leq3$, $U$ is therefore
one two-element orbit, say $\{c_0,c_1\}$, with
$e_{ij}c_k=c_i$.

Every $e_{ij}s$ lies in this unique orbit and is fixed by
$e_{ij}$, so $e_{ij}s=c_i$ for every $s\in S$. If $S=U$, its
permutation module is $Re_{00}$, via $e_{i0}\mapsto c_i$, and
is projective. Otherwise $S=\{c_0,c_1,w\}$. The map
$R\to\mathbb Q[S]$, $r\mapsto rw$, is onto and has kernel
spanned by $e_{00}-e_{01}$ and $e_{10}-e_{11}$. Thus
\[
 0\longrightarrow Re_{00}\xrightarrow{\iota}R
 \xrightarrow{r\mapsto rw}\mathbb Q[S]\longrightarrow0,
 \qquad \iota(e_{i0})=e_{i0}-e_{i1},
\]
is exact. The formula for $\iota$ is $R$-linear by the band
multiplication, and its two displayed images are independent.
Both terms on the left are projective.
\end{proof}

\begin{theorem}[Proved: a small stalk and a singleton long hom-set]
\label{CP2:A-five-middle-three-stalk-singleton-long}
Let $C$ have full linear support $0\to1\to2$, trivial endpoint
corners, middle corner $M$ of order at most five, and a singleton
hom-set $C(0,2)$. If $1\leq |C(0,1)|\leq3$, then
$BC\not\simeq\CP^2$, with no bound on $|C(1,2)|$.
The opposite statement holds as well.
\end{theorem}
\begin{proof}
Put $S=C(0,1)$, a left $M$-set, and $T=C(1,2)$, a right
$M$-set. Every corner monoid of order at most four, and every
five-element monoid other than $P_5$, is rationally acyclic;
for order five this is Lemma~\ref{CP2:A-order-five-corner}.
In these cases Theorem~\ref{CP2:A-baseline-singleton-long}
makes all positive rational cohomology products zero.
It remains to consider $M=P_5$.

Suppose first that $S$ has an $M$-fixed point $s$. Use the
support cut with first block $\{0\}$ and second block
$\{1,2\}$. Its connecting profunctor has components $S$ and
$C(0,2)$. Collapsing $S$ to $s$ and fixing the singleton second
component is an idempotent natural endomorphism: fixedness
checks naturality for $M$, and each arrow of $T$ takes every
element of $S$ to the unique long arrow. The image category
has exactly one arrow from $0$ to every object, so $0$ is an
initial object and its nerve is contractible. If $BC$ had type
$\CP^2$, Theorem~\ref{CP2:A-cut-retract} would give that same
type to the image category, a contradiction.

Now assume that $S$ has no fixed point. The preceding lemma
gives $\operatorname{pd}_R\mathbb Q[S]\leq1$. The explicit
resolution in Lemma~\ref{CP2:A-order-five-corner} gives
$\operatorname{pd}_R\mathbb Q\leq2$ for the trivial left
module; its opposite gives the same bound for the trivial
right module.

Use the actual triangular bar diagram obtained by separating
nerve strings at their strict object changes, as in
Theorem~\ref{CP2:A-baseline-singleton-long}. Its vertex spaces
are $*,BP_5,*$. Its adjacent edge spaces compute
\[
 A:\quad \mathbb Q\otimes_R^{\mathbf L}\mathbb Q[S],
 \qquad
 B:\quad \mathbb Q[T]\otimes_R^{\mathbf L}\mathbb Q,
\]
and its long-edge space is a point. Its face space computes
\[
 F:\quad \mathbb Q[T]\otimes_R^{\mathbf L}\mathbb Q[S].
\]
These descriptions retain the middle bar direction and the
actual composition map; no product model for $F$ is required.
The bounds just proved give $H_q(A;\mathbb Q)=H_q(F;\mathbb Q)=0$
for $q\geq2$, and $H_q(B;\mathbb Q)=0$ for $q\geq3$.
The vertex homology also vanishes above degree two.

The finite strict-change filtration has columns $p=0,1,2$.
Its first page is the direct sum of the rational homology of
the vertex, edge, and face bars respectively, with total
degree $p+q$. The vertex terms have total degree at most two;
the edge terms have total degree at most three; and the face
terms have total degree at most three. Hence every first-page
term of total degree at least four is zero. Convergence of
this finite filtration gives $H_4(BC;\mathbb Q)=0$, again
excluding $\CP^2$. Taking opposites proves the other orientation.
\end{proof}

\begin{corollary}[Proved: the remaining eight-arrow adjacent stalk]
\label{CP2:A-nineteen-five-middle-eight-stalk}
The nineteen-arrow linear profile $(1,5,1;3,8,1)$ is excluded.
\end{corollary}
\begin{proof}
Apply the theorem with $|S|=3$ and $|T|=8$.
\end{proof}
